\textit{Note: This section is out of syllabus for NSW HSC Mathematics Extension~1 and Extension~2. It is included for enrichment and to show one way mathematicians handle differential equations that resist standard integration.}

\subsubsection*{The basic idea}
Sometimes an ordinary differential equation has no simple closed-form solution. One strategy is to assume the solution can be written as a power series and then determine its coefficients.

For a Maclaurin series,
\[
y(x)=y(0)+y'(0)x+\frac{y''(0)}{2!}x^2+\frac{y'''(0)}{3!}x^3+\cdots
\]
the differential equation itself can be used to compute successive derivatives at $x=0$.

\subsubsection*{Example 1: Solve $\boldsymbol{y'=y,\; y(0)=1}$}
From the differential equation,
\[
y'(0)=y(0)=1.
\]
Differentiating again gives $y''=y'$, so $y''(0)=1$, and similarly every higher derivative at $0$ is also $1$.

Therefore,
\[
y(x)=1+x+\frac{x^2}{2!}+\frac{x^3}{3!}+\cdots=e^x.
\]

\subsubsection*{Example 2: Solve $\boldsymbol{y'=x+y,\; y(0)=1}$}
We have
\[
y'(0)=0+y(0)=1.
\]
Differentiate the equation:
\[
y''=1+y',
\]
so
\[
y''(0)=2.
\]
Differentiate again:
\[
y'''=y'',
\]
so $y'''(0)=2$, and the same pattern continues.

This gives the beginning of the series
\[
y(x)=1+x+x^2+\frac{x^3}{3}+\frac{x^4}{12}+\cdots
\]
which matches the exact solution, $y(x)=2e^x-x-1$, as can be verified by substitution.

\subsubsection*{Takeaway}
Power-series methods turn a differential equation into a coefficient-finding problem. They are one of the main tools used when exact antiderivatives are not enough.
